% Chapter 3 of The Replica Rosetta Stone: the Sherrington--Kirkpatrick model.
\rstchapter{The Sherrington--Kirkpatrick Model}{Replica symmetry and its breaking}{broken replicas\dict interpolation}{ch:sk}
\stonefoot{proved}
\begin{refsection}[chapters/03-sk.bib]
\newcommand{\skRS}{\phi_{\mathrm{RS}}}
\newcommand{\sksech}{\mathop{\mathrm{sech}}\nolimits}
\newcommand{\skAT}{\mathrm{AT}}
\newcommand{\skPar}{\mathcal{P}}

\begin{chaptersummary}
Feed pure noise to the posterior of $\Z_2$ synchronization and you obtain the Sherrington--Kirkpatrick model, whose free entropy physicists computed with replicas: replica symmetry, its instability below the de~Almeida--Thouless line, and Parisi's formula. With Gaussian interpolation in place of replicas, we prove that the limit exists, that the replica-symmetric formula is an upper bound with an exact remainder at every size, and that it is exact at high temperature without a field, by the physicists' own computation at $n=2$. Exact enumeration checks the remainder identity on both sides for up to $14$ spins, and shows at $16$ spins why small systems cannot see the breaking of replica symmetry. The Parisi formula is a theorem, and so is full replica-symmetry breaking just below the critical temperature; replica symmetry up to the line with a field, full replica-symmetry breaking in the rest of the low-temperature phase, and the overlap law were proved in 2026 preprints, not yet refereed.
\end{chaptersummary}

\section{The question: a posterior that finds structure in noise}\label{sec:sk:question}

In $\Z_2$ synchronization a statistician observes $A=\frac\lambda Nxx^\top+W$, with $x\in\{\pm1\}^N$ unknown and $W$ a GOE matrix (\cref{ch:spiked}). With independent priors $\Prob(x_i=\pm1)=e^{\pm h}/(2\cosh h)$, the posterior computed from the entries above the diagonal (the diagonal carries no information, since $x_i^2=1$) is
\[
\Prob(x\mid A)\propto\exp\Bigl(\lambda\sum_{i<j}A_{ij}x_ix_j+h\sum_{i=1}^Nx_i\Bigr).
\]
Suppose that there is no signal, $A=W$, so that $A_{ij}=g_{ij}/\sqrt N$ for $i<j$ with independent standard Gaussians $g_{ij}$, and that the statistician computes this posterior anyway. The result is the Gibbs measure of the \term{Sherrington--Kirkpatrick model}\index{Sherrington--Kirkpatrick model} \parencite{SherringtonKirkpatrick1975} at inverse temperature $\beta=\lambda$ and field $h$: on $\Sigma_N=\{-1,+1\}^N$,
\begin{equation}\label{eq:sk:hamiltonian}
H_N(\sigma)=X_N(\sigma)+h\sum_{i=1}^N\sigma_i,\qquad X_N(\sigma)=\frac\beta{\sqrt N}\sum_{1\le i<j\le N}g_{ij}\sigma_i\sigma_j,
\end{equation}
with Gibbs weights $e^{H_N(\sigma)}/Z_N$. The \term{free entropy}\index{free entropy!of SK} is $\phi_N(\beta,h)=\frac1N\E\log Z_N$, with $\E$ the average over the couplings; the physicist's free energy per spin is $-\phi_N/\beta$, at temperature $T=1/\beta$ and physical field $H=h/\beta$.\sidenote{Up to explicit factors $Z_N$ is the likelihood ratio of the spiked model against pure noise: with the entries above the diagonal observed, the evidence ratio is $(2\cosh h)^{-N}e^{-\beta^2(N-1)/4}Z_N$. So $\phi_N$ is a log-evidence per variable, and the formula for $\E Z_N$ below says that a likelihood ratio has mean one under the null.} \Cref{ch:nishimori} treats the companion problem in which the spike is present and $\lambda$ is known; there replica symmetry is a theorem at every $\lambda$.

We write $\gibbs{\cdot}$ for the Gibbs average and use \term{replicas}\index{replicas}: independent samples $\sigma^1,\sigma^2,\dots$ from the same Gibbs measure, the couplings fixed. The \term{overlap}\index{overlap} of two replicas is $R_{12}=\frac1N\sum_i\sigma^1_i\sigma^2_i$, and $R_{\sigma\tau}=\frac1N\ip\sigma\tau$; a probabilist would call $\E\gibbs{R_{12}^2}$ the mean squared normalized inner product of two independent posterior samples.\sidenote{At $h=0$ the Gibbs measure is invariant under $\sigma\mapsto-\sigma$, so $\gibbs{\sigma_i}=0$ and the law of $R_{12}$ is symmetric; zero-field statements therefore concern $\abs{R_{12}}$, or a symmetrized measure.} One covariance computation drives the chapter: since $\sum_{i<j}a_ia_j=\frac12((\sum_ia_i)^2-\sum_ia_i^2)$, taking $a_i=\sigma_i\tau_i$ gives
\begin{equation}\label{eq:sk:cov}
\E X_N(\sigma)X_N(\tau)=\frac{\beta^2}N\sum_{i<j}\sigma_i\sigma_j\tau_i\tau_j=\frac{N\beta^2}2\Bigl(R_{\sigma\tau}^2-\frac1N\Bigr),
\end{equation}
a covariance $N\xi(R_{\sigma\tau})+O(1)$ with $\xi(q)=\beta^2q^2/2$. In particular $\E e^{X_N(\sigma)}=e^{\beta^2(N-1)/4}$, and Jensen's inequality gives the \term{annealed bound}\index{annealed bound}
\begin{equation}\label{eq:sk:annealed}
\phi_N(\beta,h)\le\frac1N\log\E Z_N=\log(2\cosh h)+\frac{\beta^2(N-1)}{4N}.
\end{equation}
At small $\beta$ this bound should be sharp and two posterior samples nearly orthogonal; at large $\beta$ the posterior should concentrate on many configurations that fit the noise and lie far apart.

\begin{physicsbox}[Sherrington--Kirkpatrick 1975; de~Almeida--Thouless 1978; Parisi 1979--80]
Let $q$ solve $q=\E\tanh^2(\beta\sqrt q\,Z+h)$, taking the positive solution when $h=0$ and $\beta>1$. If
\[
\skAT(\beta,h)\defeq\beta^2\,\E\sksech^4(\beta\sqrt q\,Z+h)\le1,
\]
the de~Almeida--Thouless condition, then the overlap concentrates at $q$ and the free entropy has the replica-symmetric value
\[
\phi(\beta,h)=\skRS(\beta,h;q)=\log2+\E\log\cosh(\beta\sqrt q\,Z+h)+\frac{\beta^2}4(1-q)^2 .
\]
Otherwise replica symmetry is broken: $\phi$ is the minimum of Parisi's functional \eqref{eq:sk:parisi} over probability measures $\mu$ on $[0,1]$, the minimizer is not a point mass, the law of $R_{12}$ converges to it (symmetrized when $h=0$), and the overlaps of many replicas are ultrametric. At $h=0$ the transition is at $\beta=1$, and the ground-state energy per spin is $-0.7632$; replica symmetry predicts $-\sqrt{2/\pi}=-0.7979$, with entropy $-1/(2\pi)$.
\end{physicsbox}

How do physicists derive this, how much of the derivation survives, and what is proved?

\section{The physics: replicas, their symmetry, and its breaking}\label{sec:sk:physics}

\subsection{The replicated partition function}\label{sec:sk:replicas}

The \term{replica trick}\index{replica trick} (\cref{ch:freeentropy}) computes $\E Z_N^n$ for integer $n$, a sum over $n$ copies of the system, and recovers $\E\log Z_N$ from $\E Z_N^n=1+n\E\log Z_N+O(n^2)$ as $n\to0$.

\begin{computation}{The replicated partition function}\label{comp:sk:replicated}
\step \emph{Average one Gaussian.} For integer $n\ge1$, $\E Z_N^n=\sum_{\sigma^1,\dots,\sigma^n}e^{h\sum_{a,i}\sigma^a_i}\,\E\exp\bigl(\sum_aX_N(\sigma^a)\bigr)$, and by \eqref{eq:sk:cov} the Gaussian $\sum_aX_N(\sigma^a)$ has variance $n\beta^2(N-1)/2+\sum_{a\ne b}\bigl(\frac{N\beta^2}2R_{ab}^2-\frac{\beta^2}2\bigr)$, with $R_{ab}=R_{\sigma^a\sigma^b}$. Since $\E e^G=e^{\Var(G)/2}$,
\begin{equation}\label{eq:sk:replicated}
\E Z_N^n=e^{nN\beta^2/4-n^2\beta^2/4}\sum_{\sigma^1,\dots,\sigma^n}\exp\Bigl(\frac{N\beta^2}2\sum_{a<b}R_{ab}^2+h\sum_{a,i}\sigma^a_i\Bigr).
\end{equation}
\step \emph{Linearize the squares.} The Hubbard--Stratonovich identity $e^{N\beta^2R^2/2}=\sqrt{N\beta^2/(2\pi)}\int_\R e^{-N\beta^2Q^2/2+N\beta^2QR}\dd Q$, with one variable $Q_{ab}$ for each pair $a<b$, makes the exponent linear in the overlaps, and the spin sum factorizes over sites:
\[
\E Z_N^n=e^{nN\beta^2/4-n^2\beta^2/4}\Bigl(\frac{N\beta^2}{2\pi}\Bigr)^{n(n-1)/4}\int e^{NS_n(Q)}\prod_{a<b}\dd Q_{ab},
\]
\[
S_n(Q)=-\frac{\beta^2}2\sum_{a<b}Q_{ab}^2+\log\sum_{s\in\{\pm1\}^n}\exp\Bigl(\beta^2\sum_{a<b}Q_{ab}s_as_b+h\sum_as_a\Bigr).
\]
\step \emph{Laplace's method at fixed $n$.} Since $S_n$ is smooth and tends to $-\infty$ quadratically, $\frac1N\log\E Z_N^n\to\frac{n\beta^2}4+\max_QS_n(Q)$ (\cref{thm:toolkit:laplace}). For each fixed integer $n$ this step is a theorem.
\end{computation}

Two uncontrolled steps remain: the exchange of $N\to\infty$ with $n\to0$, and the continuation of $n\mapsto\max_QS_n(Q)$ from the integers to $n$ near $0$, which the integer moments alone cannot justify (\cref{ch:freeentropy}). With $Q^\star$ a stationary point of $S_n$ chosen by an ansatz, the prediction is $\phi=\lim_{n\to0}\frac1n\bigl(\frac{n\beta^2}4+S_n(Q^\star)\bigr)$.

\subsection{Replica symmetry}\label{sec:sk:rs}

\begin{heuristic}[The replica-symmetric solution]
Take $Q_{ab}=q$ for all $a<b$. Then $\sum_{a<b}Q_{ab}s_as_b=\frac q2(\sum_as_a)^2-\frac{nq}2$, and a second Hubbard--Stratonovich variable $z$ decouples the replicas:
\[
S_n(q)=-\frac{\beta^2}2\,\frac{n(n-1)}2\,q^2-\frac{n\beta^2q}2+\log\E_z\bigl(2\cosh(\beta\sqrt q\,z+h)\bigr)^n.
\]
Dividing by $n$ and letting $n\to0$, with $\frac1n\log\E_zY^n\to\E_z\log Y$, gives
\begin{equation}\label{eq:sk:phirs}
\skRS(\beta,h;q)=\log2+\E\log\cosh(\beta\sqrt q\,Z+h)+\frac{\beta^2}4(1-q)^2.
\end{equation}
Gaussian integration by parts, $\E[Z\tanh(aZ+h)]=a\E\sksech^2(aZ+h)$, gives
\begin{equation}\label{eq:sk:fixedpoint}
\partial_q\skRS(\beta,h;q)=\frac{\beta^2}2\bigl(q-\E\tanh^2(\beta\sqrt q\,Z+h)\bigr),
\end{equation}
so the order parameter solves $q=\E\tanh^2(\beta\sqrt q\,Z+h)$: it is the overlap $\gibbs{s_as_b}$ of two replicas under the replicated single-site measure.
\end{heuristic}

The continuation has a cost. For integer $n\ge2$ the saddle point maximizes $S_n$, but the coefficient $-\frac{\beta^2}4(n-1)$ of $q^2$ in $S_n(q)/n$ changes sign at $n=1$, and as $n\to0$ the physical $q$ \emph{minimizes} $\skRS$. At $h=0$ the only solution for $\beta\le1$ is $q=0$, with the annealed value; for $\beta>1$ a positive solution appears, $q=0$ becomes a local maximum, and the physicist takes the minimum (\cref{rem:sk:fixedpoint}).

\begin{computation}{Replica symmetry at zero temperature}\label{comp:sk:zeroT}
\step \emph{The order parameter.} At $h=0$ and $\beta\to\infty$, $q\to1$, and substituting $u=\beta\sqrt q\,z$ in $1-q=\E\sksech^2(\beta\sqrt q\,Z)$, with $\int_\R\sksech^2u\dd u=2$, gives $\beta(1-q)\to2/\sqrt{2\pi}=\sqrt{2/\pi}$.
\step \emph{Energy.} Since $\log\cosh x=\abs x-\log2+\log(1+e^{-2\abs x})$, $\skRS=\beta\sqrt q\,\E\abs Z+\E\log(1+e^{-2\beta\sqrt q\abs Z})+\frac14(\beta(1-q))^2$, and the middle term is $O(1/\beta)$. As $\phi=-\beta f=-\beta e+s$, the ground-state energy is $e_0=-\lim\skRS/\beta=-\E\abs Z=-\sqrt{2/\pi}=-0.7979$.
\step \emph{Entropy.} With $\beta(1-\sqrt q)=\beta(1-q)/(1+\sqrt q)\to\frac12\sqrt{2/\pi}$,
\begin{align*}
s_0=\lim\bigl(\skRS+\beta e_0\bigr)&=\lim\Bigl(-\sqrt{\tfrac2\pi}\,\beta(1-\sqrt q)+\tfrac14\bigl(\beta(1-q)\bigr)^2\Bigr)\\
&=-\frac1\pi+\frac1{2\pi}=-\frac1{2\pi}.
\end{align*}
\end{computation}

Ising spins have a nonnegative entropy,\sidenote{At $h=0$, $\phi_N-\beta\partial_\beta\phi_N=\frac1N\E\bigl[\log Z_N-\gibbs{H_N}\bigr]$ is the mean Shannon entropy of the Gibbs measure divided by $N$, which lies in $[0,\log2]$.} so $s_0=-0.159$ is impossible: replica symmetry fails at low temperature. The script \texttt{03-sk-replica.py} confirms the limits ($\beta(1-q)=0.79820$, $e=-0.797884$ and $s=-0.15863$ at $T=0.001$) and finds that the replica-symmetric entropy changes sign at $T=0.2716$ (\cref{fig:sk:phase}(b)).

\subsection{The de~Almeida--Thouless instability}\label{sec:sk:at}

\begin{heuristic}[The replicon]
The Hessian of $S_n$ at $Q_{ab}=q$ is $M=-\beta^2I+\beta^4K$, with $K_{(ab),(cd)}=\gibbs{s_as_bs_cs_d}_n-\gibbs{s_as_b}_n\gibbs{s_cs_d}_n$ for the replicated single-site measure $\gibbs\cdot_n$. As $n\to0$ its entries are $1-q^2$ for equal pairs, $q-q^2$ for pairs sharing one index and $r-q^2$ for disjoint pairs, with $r=\E\tanh^4(\beta\sqrt q\,Z+h)$. A matrix indexed by pairs with three such entries, $c_2$, $c_1$ and $c_0$ for pairs sharing two, one or no indices, acts on the space of $\eta$ with $\sum_b\eta_{ab}=0$ for every $a$, of dimension $n(n-3)/2$, as multiplication by $c_2-2c_1+c_0$ (\cref{ex:sk:replicon}). Since $1-2q+r=\E(1-\tanh^2)^2$, this \term{replicon}\index{replicon} eigenvalue is
\begin{equation}\label{eq:sk:replicon}
\lambda_{\mathrm R}=-\beta^2+\beta^4\bigl[(1-q^2)-2(q-q^2)+(r-q^2)\bigr]=-\beta^2\bigl(1-\skAT(\beta,h)\bigr).
\end{equation}
At high temperature $\lambda_{\mathrm R}\approx-\beta^2<0$. De~Almeida and Thouless \parencite*{deAlmeidaThouless1978} took its vanishing as the limit of stability of replica symmetry, which is stable when $\skAT(\beta,h)\le1$.
\end{heuristic}

The argument is heuristic twice over: the multiplicity $n(n-3)/2$ tends to $-1$, and after the inversion at $n=1$ no rule fixes the sign each eigenvalue should have. The \term{de~Almeida--Thouless line}\index{de Almeida--Thouless line} is nevertheless sharp: at $h=0$ it is the point $\beta=1$, and with a field it starts at $(T,H)=(1,0)$ with $H^2\simeq\frac43(1-T)^3$ (\cref{ex:sk:smallfield}) and rises as the temperature falls (\cref{fig:sk:phase}(a)).

\begin{figure}[t]
\begin{fullwidth}
\centering
\begin{tikzpicture}
\begin{groupplot}[group style={group size=2 by 1,horizontal sep=17mm},rplot,width=80mm,height=44mm]
\nextgroupplot[xlabel={temperature $T=1/\beta$},ylabel={field $H=h/\beta$},xmin=0,xmax=1.3,ymin=0,ymax=2.3,
  xtick={0,0.25,0.5,0.75,1,1.25},ytick={0,0.5,1,1.5,2},title={(a) the de~Almeida--Thouless line},legend pos=north east]
\addplot[physpred] table[x=T,y=H] {figures/data/03-sk-atline.dat};
\addplot[guide] table[x=T,y=Hsmall] {figures/data/03-sk-atline.dat};
\legend{{$\skAT(\beta,h)=1$},{$H^2=\frac43(1-T)^3$}}
\node[font=\footnotesize,text=ink] at (axis cs:0.95,1.15) {replica symmetric};
\node[font=\footnotesize,text=ink] at (axis cs:0.16,0.25) {broken};
\nextgroupplot[xlabel={temperature $T$},ylabel={entropy per spin at $h=0$},xmin=0,xmax=1.5,ymin=-0.2,ymax=0.65,
  xtick={0,0.25,0.5,0.75,1,1.25,1.5},title={(b) the entropy of the ans\"atze},legend pos=south east]
\addplot[physpred] table[x=T,y=sRS] {figures/data/03-sk-rs.dat};
\addplot[physpred2] table[x=T,y=s1] {figures/data/03-sk-rs.dat};
\addplot[guide] coordinates {(0,-0.159155) (1.5,-0.159155)};
\addplot[muted,line width=0.4pt,forget plot] coordinates {(0,0) (1.5,0)};
\legend{replica symmetric,one-step RSB,$-1/(2\pi)$}
\end{groupplot}
\end{tikzpicture}
\end{fullwidth}
\caption{The replica predictions. (a) The de~Almeida--Thouless line in the plane of temperature and physical field, from \eqref{eq:sk:fixedpoint} and $\skAT=1$ at $65$ temperatures; dashed, its small-field law. Replica symmetry holds above the line (proved at $h=0$ by \textcite{AizenmanLebowitzRuelle1987}, for $h>0$ in the 2026 preprints of \cref{sec:sk:beyond}) and fails below it \parencite{Toninelli2002}. (b) The entropy $s=\phi-\beta\partial_\beta\phi$ of the replica-symmetric and one-step solutions at $h=0$, the latter minimized over $(q,m)$ at $38$ temperatures below $T=1$. The replica-symmetric entropy is negative below $T=0.2716$ and tends to $-1/(2\pi)$ (\cref{comp:sk:zeroT}); one step of breaking reduces the defect without removing it ($-0.023$ at $T=0.05$). Script: \texttt{figures/scripts/03-sk-replica.py}.}\label{fig:sk:phase}
\end{figure}

\subsection{Parisi's hierarchy}\label{sec:sk:parisi}

\begin{heuristic}[Breaking replica symmetry in steps]
\textcite{Parisi1979b} divided the $n$ replicas into $n/m$ groups of $m$, with $Q_{ab}=q_1$ for $a\ne b$ in the same group and $q_0$ otherwise. One Hubbard--Stratonovich variable $z_1$ per group and a common one $z_0$ decouple the replicas, with $\Xi=h+\beta\sqrt{q_0}\,z_0+\beta\sqrt{q_1-q_0}\,z_1$ as the local field. For integer $n$ the group size satisfies $1\le m\le n$; Parisi continued it to $0\le m\le1$, the interval between the two ends once $n\to0$, and obtained
\begin{equation}\label{eq:sk:p1rsb}
\skPar_1(q_0,q_1,m)=\log2+\frac1m\E_{z_0}\log\E_{z_1}\cosh^m\Xi+\frac{\beta^2}4\bigl[(1-q_1)^2-m(q_1^2-q_0^2)\bigr].
\end{equation}
Iterating $k$ times gives nested groups, overlaps $q_0<\dots<q_k$ and sizes $0\le m_1\le\dots\le m_k\le1$, all encoded in the function $x(s)=\mu([0,s])$ of a probability measure $\mu$ on $[0,1]$, with $x=m_p$ between $q_{p-1}$ and $q_p$ \parencite{Parisi1980b}. As $k\to\infty$ the free entropy becomes the minimum over $\mu$ of
\begin{equation}\label{eq:sk:parisi}
\skPar_{\beta,h}(\mu)=\log2+u_\mu(0,h)-\frac{\beta^2}2\int_0^1s\,\mu([0,s])\dd s,
\end{equation}
where $u_\mu$ solves the \term{Parisi equation}\index{Parisi formula!equation}
\begin{equation}\label{eq:sk:pde}
\partial_su+\frac{\beta^2}2\Bigl(\partial_y^2u+\mu([0,s])\,(\partial_yu)^2\Bigr)=0\quad(0\le s<1),\qquad u(1,y)=\log\cosh y.
\end{equation}
Where $x\equiv m>0$ on $[a,b]$, the substitution $e^{mu}$ gives the backward heat equation, so $e^{mu(a,y)}=\E e^{mu(b,y+\beta\sqrt{b-a}\,Z)}$; where $x\equiv0$, $u(a,y)=\E u(b,y+\beta\sqrt{b-a}\,Z)$. These rules give \eqref{eq:sk:phirs} for $\mu=\delta_q$ and \eqref{eq:sk:p1rsb} for $\mu=m\delta_{q_0}+(1-m)\delta_{q_1}$ (\cref{ex:sk:pde}).
\end{heuristic}

\textcite{Parisi1983} identified $x(q)$ with the probability that two replicas have overlap at most $q$. Below the transition the Gibbs measure splits into many \term{pure states}\index{pure states},\sidenote{In the physical picture their weights are those of a Poisson--Dirichlet point process, the law of the Gibbs weights of the random energy model at low temperature (\cref{ch:rem}).} whose overlaps, \textcite{MezardEtAl1984a} found, are \term{ultrametric}\index{ultrametricity}: for three replicas, the two smallest of $R_{12}$, $R_{13}$, $R_{23}$ are equal. In numerical solutions at $h=0$ below $T=1$ \parencite{CrisantiRizzo2002}, $\mu$ has a continuous part on $[0,q_{\max}]$ and an atom at $q_{\max}$: \term{full replica-symmetry breaking}\index{replica-symmetry breaking!full}.

Each step lowers the value: at $h=0$ and $T=0.5$, $\skRS=1.584786$ and $\min\skPar_1=1.582610$ (at $q_0=0$, $q_1=0.60132$, $m=0.20424$). For the ground state, the zero-temperature version of \eqref{eq:sk:parisi} \parencite{AuffingerChen2017} with $k=0,1,2$ steps gives $0.797885$, $0.765259$ and $0.763608$, against $0.7631667265(6)$ extrapolated from $42$ steps by \textcite{OppermannSchmidtSherrington2007}; \textcite{Parisi1980b} had $0.7633\pm0.0001$, as reported by \textcite{CrisantiRizzo2002}. In this derivation the Hubbard--Stratonovich and Laplace steps at fixed $n$ are rigorous; the exchange of limits, the continuation in $n$, the ansatz, the inverted rule that picks the stationary point and the stability analysis are not.

\section{The stone}\label{sec:sk:stone}

\begin{stone}{the replica computation}{Gaussian interpolation}
\stonestep{The thermodynamic limit}{The free entropy is assumed to self-average and to have a limit.}{Gaussian concentration, and an interpolation between one system and two independent halves that makes $N\phi_N$ superadditive (\cref{thm:sk:guerratoninelli}). An idea the physics does not need.}
\stonestep{Replicate}{Compute $\E Z_N^n$ for integer $n$ (\cref{comp:sk:replicated}); the overlap matrix $Q_{ab}$ is the order parameter; Laplace in $N$ at fixed $n$; then $n\to0$.}{No moment of $Z_N$ is computed. Interpolate between $H_N$ and a one-body field $\beta\sqrt q\sum_iz_i\sigma_i$; integration by parts turns $\frac{\dd}{\dd t}\E\log Z_t$ into averages over \emph{two} replicas (\cref{lem:sk:interp}).}
\stonestep{The replica-symmetric saddle}{$Q_{ab}=q$; stationarity gives $q=\E\tanh^2(\beta\sqrt q\,Z+h)$; since $n(n-1)/2<0$ for $0<n<1$, the maximum over $Q$ becomes a minimum over $q$.}{Guerra's bound $\phi_N\le\skRS(q)$ holds for every $q$ and $N$, with an exact remainder (\cref{thm:sk:rsbound}). No counterpart for the continuation in $n$: the minimum over $q$ comes from the sign of the remainder.}
\stonestep{Two replicas at high temperature}{$\E Z_N^2$ by Hubbard--Stratonovich: the Gaussian integral around $Q_{12}=0$ converges when $\beta<1$, so quenched equals annealed.}{The same identity bounds $\E Z_N^2/(\E Z_N)^2$ at every $N$ (\cref{lem:sk:secondmoment}), and with concentration proves $\phi=\log2+\beta^2/4$ for $\beta<1$ (\cref{thm:sk:hightemp}). The physical computation is the proof.}
\stonestep{Stability: the replicon}{The replicon eigenvalue $-\beta^2(1-\skAT)$, of multiplicity $n(n-3)/2\to-1$; replica symmetry is stable if and only if $\skAT\le1$.}{The remainder $\frac{\beta^2}4\int_0^1\E\gibbs{(R_{12}-q)^2}_t\dd t$ tends to $0$ if and only if the overlap concentrates. A one-step perturbation of Guerra's bound has $m$-derivative $\frac{\beta^2\delta^2}4(\skAT-1)+O(\delta^3)$, positive below the line (\cref{thm:sk:toninelli}); above it, 2026 preprints prove symmetry (\cref{thm:sk:atline}).}
\stonestep{Parisi's hierarchy}{Nested groups of replicas, sizes $0\le m_1\le\dots\le1$ as $n\to0$; the function $x(q)$ and the Parisi equation; Parisi's solution is taken to be the right stationary point.}{Guerra replaces the hierarchy by nested averages over one-body fields (Ruelle probability cascades, in later language): again a mean-square remainder, so $\phi_N\le\skPar(\mu)$ for every $\mu$. The lower bound needs a new idea: Talagrand's two coupled replicas, or Aizenman--Sims--Starr with ultrametricity (\cref{thm:sk:parisi}).}
\stonestep{The meaning of $x(q)$}{$x$ is the distribution function of $R_{12}$; pure states carry Poisson--Dirichlet weights and ultrametric overlaps.}{Ghirlanda--Guerra identities, from the concentration of the energy under small perturbations, imply ultrametricity; for SK itself, a 2026 preprint identifies the limit of the overlap array (\cref{thm:sk:overlap}).}
\end{stone}

Read across, the stone shows where the routes meet and part. The proof justifies the replica functionals as upper bounds at every $N$, and at $n=2$ the Hubbard--Stratonovich computation is literally a proof. It replaces the rest: the continuation in $n$ by an interpolation in a real parameter, the inverted saddle-point rule by the sign of a remainder, the stability analysis by a perturbation of a bound. And it needs ideas the physics never needed, superadditivity, cascades and ultrametricity. The formula and its instability line date from 1975--78; the matching rigorous statements, from 1987 at zero field, 2002 below the line, and 2026, in preprints, above it.

\section{The proof: interpolation, from the limit to high temperature}\label{sec:sk:proof}

The first two theorems of this section rest on one lemma. Let $\Sigma$ be a finite set with positive weights $w(\sigma)$, and let $(X_\sigma)$ and $(Y_\sigma)$ be independent centered Gaussian processes indexed by $\Sigma$, with covariances $C^X$ and $C^Y$. For $t\in[0,1]$ let $\gibbs\cdot_t$, extended to replicas, be the Gibbs average of
\[
Z_t=\sum_{\sigma\in\Sigma}w(\sigma)\exp\bigl(\sqrt t\,X_\sigma+\sqrt{1-t}\,Y_\sigma\bigr),\qquad\text{and let }\psi(t)=\E\log Z_t.
\]

\begin{lemma}[Interpolation of free entropies]\label[lemma]{lem:sk:interp}
The function $\psi$ is continuously differentiable on $[0,1]$, and
\begin{equation}\label{eq:sk:interp}
\psi'(t)=\frac12\E\Bigl\langle C^X(\sigma,\sigma)-C^Y(\sigma,\sigma)\Bigr\rangle_t-\frac12\E\Bigl\langle C^X(\sigma^1,\sigma^2)-C^Y(\sigma^1,\sigma^2)\Bigr\rangle_t.
\end{equation}
\end{lemma}

\begin{proof}
The plan is to apply Gaussian interpolation (\cref{lem:toolkit:interpolation}) to a log-partition function. Put $W_t=\sqrt tX+\sqrt{1-t}Y$. That lemma asks for a twice continuously differentiable $F$ on $\R^\Sigma$ such that $F$ and its partial derivatives of first and second order have moderate growth (for every $a>0$, each is at most a constant times $e^{a\abs x^2}$); it then says that $t\mapsto\E F(W_t)$ is continuously differentiable on $[0,1]$, with
\[
\frac{\dd}{\dd t}\E F(W_t)=\frac12\sum_{\sigma,\tau}\bigl(C^X-C^Y\bigr)(\sigma,\tau)\,\E\,\partial_\sigma\partial_\tau F(W_t).
\]
Take the infinitely differentiable function $F(x)=\log\sum_\sigma w(\sigma)e^{x_\sigma}$ and put $G_\sigma=w(\sigma)e^{x_\sigma}/\sum_\tau w(\tau)e^{x_\tau}\in[0,1]$. Then $\abs{F(x)}\le\abs{\log\sum_\sigma w(\sigma)}+\max_\sigma\abs{x_\sigma}$,
\[
\partial_\sigma F=G_\sigma,\qquad\partial_\sigma\partial_\tau F=\ind\{\sigma=\tau\}G_\sigma-G_\sigma G_\tau ,
\]
so $F$ grows at most linearly and its first and second derivatives are bounded by $1$; all three have moderate growth, and the lemma applies. At $x=W_t$ the $G_\sigma$ are the Gibbs weights of $Z_t$, so the sum is $\gibbs{C(\sigma,\sigma)}_t-\gibbs{C(\sigma^1,\sigma^2)}_t$ with $C=C^X-C^Y$: this is \eqref{eq:sk:interp}, which the appendix also records as \cref{cor:toolkit:gibbs}. Its right side is bounded by $\max\abs{C^X}+\max\abs{C^Y}$, so $\psi$ is Lipschitz and $\psi(1)-\psi(0)=\int_0^1\psi'(t)\dd t$.
\end{proof}

The derivative sees only the covariances, on the diagonal and between two replicas.\sidenote{With $\log Z_t$ replaced by a smoothed maximum, the same computation gives the comparison inequalities of Slepian and Gordon (\cref{thm:toolkit:slepian}; \cref{ch:gordon}).}

\subsection{The thermodynamic limit}\label{sec:sk:limit}

\begin{theorem}[Guerra--Toninelli]\label{thm:sk:guerratoninelli}
For every $\beta\ge0$, $h\in\R$ and $N_1,N_2\ge1$,
\[
(N_1+N_2)\,\phi_{N_1+N_2}(\beta,h)\ge N_1\phi_{N_1}(\beta,h)+N_2\phi_{N_2}(\beta,h).
\]
Consequently the limit $\phi(\beta,h)=\lim_N\phi_N(\beta,h)$ exists, equals $\sup_N\phi_N(\beta,h)$, and is at most $\log(2\cosh h)+\beta^2/4$. If the couplings of all sizes are taken from one array $(g_{ij})_{1\le i<j}$, then $\frac1N\log Z_N\to\phi(\beta,h)$ almost surely.
\end{theorem}

\begin{proof}
The plan is to interpolate between the system of size $N=N_1+N_2$ and two independent systems on its first $N_1$ and last $N_2$ spins, and to read the sign of the derivative from the convexity of $q\mapsto q^2$. Write $\sigma=(\sigma^{(1)},\sigma^{(2)})$, let $X^{(1)}$ and $X^{(2)}$ be the disorder parts \eqref{eq:sk:hamiltonian} of two independent systems of sizes $N_1$ and $N_2$, independent of $X_N$, and apply \cref{lem:sk:interp} with $\Sigma=\Sigma_N$, $w(\sigma)=e^{h\sum_i\sigma_i}$, $X=X_N$ and $Y_\sigma=X^{(1)}(\sigma^{(1)})+X^{(2)}(\sigma^{(2)})$. Then $\psi(1)=N\phi_N$, and at $t=0$ the partition function factorizes, so $\psi(0)=N_1\phi_{N_1}+N_2\phi_{N_2}$.

With $R_1$ and $R_2$ the overlaps of the two blocks, so that $NR=N_1R_1+N_2R_2$, \eqref{eq:sk:cov} gives
\begin{align*}
C^X(\sigma,\sigma)-C^Y(\sigma,\sigma)&=\tfrac12\beta^2(N-1)-\tfrac12\beta^2(N_1-1)-\tfrac12\beta^2(N_2-1)=\tfrac12\beta^2,\\
C^X(\sigma^1,\sigma^2)-C^Y(\sigma^1,\sigma^2)&=\tfrac12\beta^2\bigl(NR^2-N_1R_1^2-N_2R_2^2\bigr)+\tfrac12\beta^2.
\end{align*}
In \eqref{eq:sk:interp} the constants cancel, and $\psi'(t)=-\frac{\beta^2}4\E\gibbs{NR^2-N_1R_1^2-N_2R_2^2}_t\ge0$, because $R=\frac{N_1}NR_1+\frac{N_2}NR_2$ and convexity give $NR^2\le N_1R_1^2+N_2R_2^2$. Hence $\psi(1)\ge\psi(0)$, the superadditivity.

Fekete's lemma turns superadditivity into convergence. Put $a_N=N\phi_N$, $a_0=0$, and fix $k\ge1$. For $N=mk+r$ with $0\le r<k$, superadditivity gives $a_N\ge ma_k+a_r$, so $a_N/N\ge\frac{mk}N\frac{a_k}k+\frac{a_r}N\to\frac{a_k}k$. Hence $\liminf_Na_N/N\ge\sup_ka_k/k\ge\limsup_Na_N/N$, and the supremum is finite, because \eqref{eq:sk:annealed} bounds every $a_k/k$ by $\log(2\cosh h)+\beta^2/4$: the limit exists, is a real number, and equals the supremum.

For the almost sure statement, $F=\log Z_N$ is a smooth function of the $N(N-1)/2$ standard Gaussians $g_{ij}$, with $\partial F/\partial g_{ij}=\frac\beta{\sqrt N}\gibbs{\sigma_i\sigma_j}$, so $\abs{\nabla F}^2\le\frac{\beta^2}N\frac{N(N-1)}2$ and $F$ is Lipschitz with constant $L=\beta\sqrt{(N-1)/2}$. Gaussian concentration (\cref{thm:toolkit:lipschitz}; \cref{thm:freeentropy:concentration}) gives, for $u>0$,
\[
\Prob\Bigl(\Bigl|\frac1N\log Z_N-\phi_N\Bigr|\ge u\Bigr)\le2\exp\Bigl(-\frac{N^2u^2}{2L^2}\Bigr)\le2e^{-Nu^2/\beta^2},
\]
summable in $N$; the Borel--Cantelli lemma, for each rational $u>0$, gives the almost sure convergence. At $\beta=0$, $\log Z_N=N\log(2\cosh h)$ is not random.
\end{proof}

Superadditivity also gives $\phi_N\le\phi$ for every $N$: finite-size measurements lie below the limit (\cref{sec:sk:measuring}).

\subsection{Guerra's replica-symmetric bound}\label{sec:sk:rsbound}

\begin{theorem}[Guerra's bound, with its remainder]\label{thm:sk:rsbound}
Let $N\ge1$, $\beta\ge0$, $h\in\R$ and $q\in[0,1]$. Let $z_1,\dots,z_N$ be independent standard Gaussians, independent of the couplings, let $\gibbs\cdot_t$ be the Gibbs average for
\[
H_{N,t}(\sigma)=\sqrt t\,X_N(\sigma)+\sqrt{1-t}\;\beta\sqrt q\sum_{i=1}^Nz_i\sigma_i+h\sum_{i=1}^N\sigma_i,
\]
extended to replicas, and let $\E$ average over the couplings and the $z_i$. Then
\begin{equation}\label{eq:sk:sumrule}
\phi_N(\beta,h)=\skRS(\beta,h;q)-\frac{\beta^2}4\int_0^1\E\bigl\langle(R_{12}-q)^2\bigr\rangle_t\dd t .
\end{equation}
In particular $\phi_N(\beta,h)\le\inf_{q\in[0,1]}\skRS(\beta,h;q)$ for every $N$, and the same holds for $\phi(\beta,h)$.
\end{theorem}

\begin{proof}
The plan is to apply \cref{lem:sk:interp} with $X=X_N$, $Y_\sigma=\beta\sqrt q\sum_iz_i\sigma_i$ and $w(\sigma)=e^{h\sum_i\sigma_i}$, and to watch the terms of order one cancel. At $t=0$ the spins decouple, $Z_0=\prod_i2\cosh(\beta\sqrt q\,z_i+h)$, and $\frac1N\psi(0)=\log2+\E\log\cosh(\beta\sqrt q\,Z+h)$; at $t=1$, $\frac1N\psi(1)=\phi_N$. The covariance of $Y$ is $\beta^2qNR_{\sigma\tau}$, that of $X$ is \eqref{eq:sk:cov}, so $C^X-C^Y$ is $\frac{\beta^2(N-1)}2-\beta^2qN$ on the diagonal and $\frac{N\beta^2}2R_{12}^2-\frac{\beta^2}2-\beta^2qNR_{12}$ between two replicas. Dividing \eqref{eq:sk:interp} by $N$, the two terms $\beta^2/(4N)$ cancel and
\[
\frac{\psi'(t)}N=\frac{\beta^2}4-\frac{\beta^2q}2-\frac{\beta^2}4\E\bigl\langle R_{12}^2-2qR_{12}\bigr\rangle_t=\frac{\beta^2}4(1-q)^2-\frac{\beta^2}4\E\bigl\langle(R_{12}-q)^2\bigr\rangle_t .
\]
Integrating from $0$ to $1$ gives \eqref{eq:sk:sumrule}. The remainder is nonnegative, which gives the bound for every $N$ and $q$, and \cref{thm:sk:guerratoninelli} passes it to the limit.
\end{proof}

The identity \eqref{eq:sk:sumrule} is Guerra's \term{sum rule}\index{sum rule}\index{Guerra's bound}, the replica-symmetric case of \textcite[Thm.~5]{Guerra2003}, exact at every $N$ in the normalization \eqref{eq:sk:hamiltonian}, which is also Guerra's. It says what replica symmetry means: the replica-symmetric formula holds in the limit if and only if $\int_0^1\E\gibbs{(R_{12}-q)^2}_t\dd t\to0$, that is, if the overlap concentrates at $q$ along the interpolation.

\begin{remark}[Which $q$]\label[remark]{rem:sk:fixedpoint}
By \eqref{eq:sk:fixedpoint} the stationary points of $q\mapsto\skRS(\beta,h;q)$ solve the fixed-point equation. For $h\ne0$ the solution $q^\star$ is unique \parencite[Prop.~1.3.8]{Talagrand2011a}, and since the derivative is negative at $q=0$ and positive at $q=1$, it is the minimizer. At $h=0$ and $\beta\le1$ the minimizer is $q=0$ (\cref{ex:sk:qzero}). At $h=0$ and $\beta>1$, $\E\tanh^2(\beta\sqrt q\,Z)=\beta^2q+O(q^2)$, so $\skRS(q)<\skRS(0)$ for small $q>0$, and \cref{thm:sk:rsbound} proves $\phi(\beta,0)<\log2+\beta^2/4$: the annealed formula fails for every $\beta>1$. Numerically the minimum is at the positive fixed point.
\end{remark}

\subsection{Replica symmetry at high temperature}\label{sec:sk:hightemp}

At $h=0$, \cref{thm:sk:rsbound} with $q=0$ is the annealed bound. Is it attained? The \term{second moment method}\index{second moment method} answers, and its computation is the replica computation at $n=2$.

\begin{lemma}[The second moment]\label[lemma]{lem:sk:secondmoment}
Let $h=0$ and $0\le\beta<1$, and let $\rho_1,\dots,\rho_N$ be independent uniform signs. For every $N\ge1$,
\[
\frac{\E Z_N^2}{(\E Z_N)^2}=e^{-\beta^2/2}\,\E\exp\Bigl(\frac{\beta^2}{2N}\Bigl(\sum_{i=1}^N\rho_i\Bigr)^2\Bigr)\le\frac{e^{-\beta^2/2}}{\sqrt{1-\beta^2}}.
\]
\end{lemma}

\begin{proof}
The plan is to compute the ratio exactly and to bound it by a Hubbard--Stratonovich linearization. By \eqref{eq:sk:cov}, $X_N(\sigma)+X_N(\tau)$ has variance $\beta^2(N-1)+N\beta^2R_{\sigma\tau}^2-\beta^2$, so $\E Z_N^2=\sum_{\sigma,\tau}\exp\bigl(\frac{\beta^2(N-1)}2+\frac{N\beta^2}2R_{\sigma\tau}^2-\frac{\beta^2}2\bigr)$, while $(\E Z_N)^2=4^Ne^{\beta^2(N-1)/2}$. For independent uniform $\sigma$ and $\tau$ the products $\sigma_i\tau_i$ are independent uniform signs with sum $NR_{\sigma\tau}$, which gives the identity. For the bound, use $e^{a^2/2}=\E e^{ag}$ with $g\sim\Normal(0,1)$ independent of the $\rho_i$, Tonelli's theorem, and $\cosh x\le e^{x^2/2}$:
\begin{align*}
\E_\rho\exp\Bigl(\frac{\beta^2}{2N}\Bigl(\sum_i\rho_i\Bigr)^2\Bigr)&=\E_g\E_\rho\exp\Bigl(\frac{\beta g}{\sqrt N}\sum_i\rho_i\Bigr)=\E_g\Bigl[\cosh\Bigl(\frac{\beta g}{\sqrt N}\Bigr)^N\Bigr]\\
&\le\E_ge^{\beta^2g^2/2}=\frac1{\sqrt{1-\beta^2}}.\qedhere
\end{align*}
\end{proof}

\begin{theorem}[Replica symmetry at high temperature without a field]\label{thm:sk:hightemp}
Let $h=0$, $0<\beta<1$ and $c_\beta=\frac14e^{\beta^2/2}\sqrt{1-\beta^2}$. For every $N\ge1$,
\[
\phi_N(\beta,0)\ge\log2+\frac{\beta^2(N-1)}{4N}-\frac{\log2}N-\beta\sqrt{\frac{\log(2/c_\beta)}N},
\]
and $\phi_N(\beta,0)\le\log2+\frac{\beta^2(N-1)}{4N}$. Consequently $\phi(b,0)=\log2+b^2/4=\inf_q\skRS(b,0;q)$ for every $b\in[0,1]$.
\end{theorem}

\begin{proof}
The plan is to show that $Z_N$ exceeds half its mean with a probability bounded below, and to let Gaussian concentration turn this into a bound on $\E\log Z_N$. The upper bound is \eqref{eq:sk:annealed}. For $N=1$ both bounds hold because $\phi_1=\log2$; let $N\ge2$, $Z=Z_N$ and $\mu=\E Z$. By the Paley--Zygmund inequality and \cref{lem:sk:secondmoment},
\[
\Prob(Z>\mu/2)\ge\frac{(1-\frac12)^2\mu^2}{\E Z^2}\ge\frac14e^{\beta^2/2}\sqrt{1-\beta^2}=c_\beta .
\]
As in the proof of \cref{thm:sk:guerratoninelli}, $\log Z$ is Lipschitz in the couplings with constant $L=\beta\sqrt{(N-1)/2}$, so $\Prob(\log Z-\E\log Z\ge s)\le e^{-s^2/(2L^2)}$ for $s>0$ (\cref{thm:toolkit:lipschitz}). Take $s=\beta\sqrt{(N-1)\log(2/c_\beta)}$, so that $e^{-s^2/(2L^2)}=c_\beta/2$. If $\E\log Z<\log(\mu/2)-s$, then
\[
\Prob(Z>\mu/2)\le\Prob(\log Z-\E\log Z>s)\le c_\beta/2<c_\beta,
\]
a contradiction. Hence $\E\log Z\ge\log\mu-\log2-s$; dividing by $N$, inserting $\frac1N\log\mu=\log2+\frac{\beta^2(N-1)}{4N}$ and using $\sqrt{N-1}\le\sqrt N$ gives the lower bound.

Both bounds converge to $\log2+\beta^2/4$. This is $\skRS(\beta,0;0)$, and also $\inf_q\skRS(\beta,0;q)$: by $\log\cosh x\ge x^2/2-x^4/12$, the difference $\skRS(\beta,0;q)-\skRS(\beta,0;0)$ is at least $\frac{\beta^2q^2}4(1-\beta^2)\ge0$ (\cref{ex:sk:qzero}). At $\beta=0$, $Z_N=2^N$ and $\phi_N=\log2$. At $\beta=1$: each $\phi_N$ is convex in $\beta$, since $\log Z_N$ is a log-sum-exp of functions linear in $\beta$, so the limit is convex, hence continuous on $(0,\infty)$, and equals $\log2+\beta^2/4$ on $(0,1)$.
\end{proof}

The rate $N^{-1/2}$ is not sharp (the true correction is of order $1/N$), but the mechanism is visible: the annealed value is exact when the second moment, the replicated partition function at $n=2$, is of the order of the first squared, and its Hubbard--Stratonovich integral converges exactly when $\beta<1$.\sidenote{The ratio $\E Z_N^2/(\E Z_N)^2$ is also the second moment of the likelihood ratio between the spiked model with $\lambda=\beta$ and pure noise. Its boundedness for $\lambda<1$ is the classical second-moment argument that the spike cannot be detected below the spectral threshold.} At $\beta>1$ the ratio grows exponentially in $N$, dominated by pairs of configurations with a macroscopic overlap.

The free entropy also controls the overlap, through an identity exact at every $N$.

\begin{lemma}[The derivative in $\beta$]\label[lemma]{lem:sk:derivative}
For every $N\ge1$, $h\in\R$ and $\beta>0$, $\partial_\beta\phi_N(\beta,h)=\frac\beta2\bigl(1-\E\gibbs{R_{12}^2}\bigr)$.
\end{lemma}

\begin{proof}
The plan is one integration by parts, together with $\sum_{i,j}\gibbs{\sigma_i\sigma_j}^2=N^2\gibbs{R_{12}^2}$, which holds because $\gibbs{\sigma^1_i\sigma^1_j\sigma^2_i\sigma^2_j}=\gibbs{\sigma_i\sigma_j}^2$. Since $\partial_\beta H_N=X_N/\beta$ and $\abs{\gibbs{X_N}}\le\max_\sigma\abs{X_N(\sigma)}$ is integrable uniformly for $\beta$ in compact sets, $\partial_\beta\phi_N=\frac1{N\beta}\E\gibbs{X_N}$. Integration by parts (\cref{thm:toolkit:gaussianibp}) gives $\E[g_{ij}\gibbs{\sigma_i\sigma_j}]=\E[\partial_{g_{ij}}\gibbs{\sigma_i\sigma_j}]=\frac\beta{\sqrt N}\bigl(1-\E\gibbs{\sigma_i\sigma_j}^2\bigr)$, so
\begin{align*}
\E\gibbs{X_N}&=\frac{\beta^2}N\sum_{i<j}\bigl(1-\E\gibbs{\sigma_i\sigma_j}^2\bigr)=\frac{\beta^2}N\Bigl(\frac{N(N-1)}2-\frac{N^2\E\gibbs{R_{12}^2}-N}2\Bigr)\\
&=\frac{N\beta^2}2\bigl(1-\E\gibbs{R_{12}^2}\bigr).\qedhere
\end{align*}
\end{proof}

With \cref{thm:sk:hightemp} and convexity, the lemma gives $\E\gibbs{R_{12}^2}=O(N^{-1/4})$ at $h=0$ and $\beta<1$ (\cref{ex:sk:overlapfromphi}): from free entropy to overlap, the stone's fifth row in its simplest case. The true order is $1/N$; the physics predicts $N\E\gibbs{R_{12}^2}\to1/(1-\beta^2)$, and \textcite{AizenmanLebowitzRuelle1987} proved a central limit theorem for $\sqrt NR_{12}$, as \textcite{KusuokaNakajima2026b} recount.

\section{The proof: Parisi's formula and the de~Almeida--Thouless line}\label{sec:sk:beyond}

The rest of the theory is stated with its sources, in this chapter's normalization (\textcite{Panchenko2013b} and \textcite{AuffingerChen2015a} write $\xi(q)=\beta^2q^2$: their $\beta$ is our $\beta/\sqrt2$).

\begin{theorem}[The Parisi formula]\label{thm:sk:parisi}
For every $\beta\ge0$ and $h\in\R$, $\lim_N\phi_N(\beta,h)=\min_\mu\skPar_{\beta,h}(\mu)$, the minimum of \eqref{eq:sk:parisi} over Borel probability measures on $[0,1]$, and the minimizer $\mu_{\beta,h}$, the \term{Parisi measure}\index{Parisi measure}, is unique.
\end{theorem}

The upper bound is Guerra's \parencite*[Thm.~3]{Guerra2003} and the lower bound Talagrand's \parencite*{Talagrand2006}, for SK and mixed even $p$-spin models; \textcite[Cor.~1]{AuffingerChen2015a} proved uniqueness. Guerra's proof runs the interpolation of \cref{thm:sk:rsbound} with a hierarchy of independent one-body fields, one for each level of a discrete $\mu$, averaged recursively with the exponents $m_p$; the remainder is again a mean square, now in auxiliary Gibbs states. In the later language of \textcite{AizenmanSimsStarr2003} and \textcite{Panchenko2013b} the interpolation runs on $\Sigma_N$ times the leaves of a Ruelle probability cascade, leaves with Poisson--Dirichlet weights whose overlaps are the $q_p$, and the remainder is $\frac{\beta^2}4\int_0^1\E\gibbs{(R_{12}-q_{\alpha^1\wedge\alpha^2})^2}_t\dd t$, with $q_{\alpha^1\wedge\alpha^2}$ the overlap of the two leaves. The lower bound needs a reason for this remainder to vanish at the optimal $\mu$: Talagrand bounded the free entropy of two coupled replicas with constrained overlap by a two-dimensional Guerra bound, and \textcite{Panchenko2014} gave a second proof, for general mixtures, from the cavity representation of \textcite{AizenmanSimsStarr2003}, the \term{Ghirlanda--Guerra identities}\index{Ghirlanda--Guerra identities} \parencite{GhirlandaGuerra1998} and ultrametricity \parencite{Panchenko2013a}. Neither idea has a counterpart in the replica computation. Below the line, Guerra's bound beats replica symmetry.

\begin{theorem}[Toninelli]\label{thm:sk:toninelli}
Let $q$ solve the fixed-point equation, taking the positive solution when $h=0$ and $\beta>1$. If $\skAT(\beta,h)>1$, then $\phi(\beta,h)<\skRS(\beta,h;q)$.
\end{theorem}

\begin{proofsketch}
Evaluate Guerra's broken bound at $\mu_m=m\delta_q+(1-m)\delta_{q+\delta}$ with $\delta>0$ small, so that $\skPar(\mu_1)=\skRS(q)$. With $a^2=\beta^2\delta$, $Y=h+\beta\sqrt q\,Z'$ and $\nu$ the law of $Z$ tilted by $\cosh(Y+aZ)$,
\[
\partial_m\skPar(\mu_m)\big|_{m=1}=\E\Bigl[\E_\nu\log\cosh(Y+aZ)-\log\cosh Y-\frac{a^2}2\Bigr]-\frac{\beta^2}4(2q\delta+\delta^2).
\]
Expanding with the heat equation, the first expectation is $\frac{a^2}2\E\tanh^2Y+\frac{a^4}4\E\sksech^4Y+O(a^6)$; the fixed-point equation cancels the first-order terms, and
\begin{equation}\label{eq:sk:toninelli}
\partial_m\skPar(\mu_m)\big|_{m=1}=\frac{\beta^2\delta^2}4\bigl(\skAT(\beta,h)-1\bigr)+O(\delta^3).
\end{equation}
If $\skAT>1$ this is positive for small $\delta$, so $\phi\le\skPar(\mu_m)<\skRS(q)$ for $m$ slightly below $1$. Omitted: the uniformity of the $O(a^6)$ term. The argument and the coefficient are \textcite{Toninelli2002}'s, who computes the second derivative in $\delta$ of the left side of \eqref{eq:sk:toninelli} at $\delta=0$, $\frac{\beta^2}2(\skAT-1)$; the route through the heat equation is ours. The script \texttt{03-sk-replica.py} confirms the coefficient: at $(\beta,h)=(2,0.5)$ the ratio of the exact derivative to it is $0.78$, $0.89$, $0.94$ and $0.97$ at $\delta=0.04$, $0.02$, $0.01$ and $0.005$, and one minus the ratio halves with $\delta$, as an $O(\delta^3)$ remainder requires.
\end{proofsketch}

The replicon eigenvalue \eqref{eq:sk:replicon} reappears in \eqref{eq:sk:toninelli}: the rigorous trace of the stability analysis. The converse remained open for $h>0$ after partial results \parencite{JagannathTobasco2017a,BrenneckeYau2022}; in February 2026 \textcite{MourratSchertzer2026}, who showed in a preprint that a generalized criterion fails for some mixed $p$-spin models, still called the SK case open. Two preprints have since claimed it.

\begin{theorem}[The de~Almeida--Thouless line with a field; preprints]\label{thm:sk:atline}
\leavevmode
\begin{enumerate}
\item \parencite[Thm.~1.1 and Prop.~1.2]{Lopatto2026a} If $\beta>0$, $h>0$ and $\skAT(\beta,h)\le1$, then $\phi(\beta,h)=\skRS(\beta,h;q)$. The set $\{\skAT=1\}$ is the graph of a smooth, strictly increasing function $h_{\mathrm{AT}}\colon(1,\infty)\to(0,\infty)$ with $h_{\mathrm{AT}}(1^+)=0$, together with the point $(1,0)$.
\item \parencite[Thm.~1.1]{KusuokaNakajima2026b} For every compact set $K\subset\{\beta>0,\ h>0,\ \skAT(\beta,h)<1\}$ there is $C_K$ such that $N\E\gibbs{(R_{12}-q)^2}_t\le C_K$ for all $t\in[0,1]$, $(\beta,h)\in K$ and $N\ge1$, and $0\le\skRS(\beta,h;q)-\phi_N(\beta,h)\le C_K/N$.
\end{enumerate}
\end{theorem}

Lopatto checks the optimality condition of \textcite{JagannathTobasco2017a} for $\mu=\delta_q$, so his proof rests on the Parisi formula; the proof of Kusuoka and Nakajima is, as they put it, self-contained, does not use the Parisi formula, and gets its rate through the sum rule \eqref{eq:sk:sumrule}. Both are unrefereed as of October 2026. Lopatto states that his paper was written with the assistance of large language models, which suggested arguments, helped draft and revise, and coded simulations. For their key step, the flatness of the one-step bound near the replica-symmetric point, Kusuoka and Nakajima \enquote{asked ChatGPT 5.6 Sol to construct a proof}; they revised the resulting argument in detail, take full responsibility for it, and used Codex to formalize the proof in Lean~4.

At $h=0$, \textcite[Thm.~1 and Cor.~2]{Zhou2026} proved, in a refereed paper, that for $\beta>1$ close to $1$ the Parisi measure is supported on an interval $[0,q_\beta]$, with a smooth density on $[0,q_\beta)$ and a single atom, at $q_\beta$: full replica-symmetry breaking just below the critical temperature. The rest is in preprints.

\begin{theorem}[Full replica-symmetry breaking and the overlap law; preprints]\label{thm:sk:overlap}
\leavevmode
\begin{enumerate}
\item \parencite[Thm.~1.1]{Lopatto2026b} For every $\beta>1$ and $h=0$ there are $q_\beta,c_\beta\in(0,1)$ and a nonnegative $\rho_\beta$, smooth up to the endpoints, with $\mu_{\beta,0}(\dd s)=\rho_\beta(s)\ind_{[0,q_\beta)}(s)\dd s+c_\beta\delta_{q_\beta}(\dd s)$ and $\supp\mu_{\beta,0}=[0,q_\beta]$.
\item \parencite[Thm.~1.1]{Lopatto2026c} For every $\beta>0$ and $h>0$ the support of $\mu_{\beta,h}$ is an interval $[q_-,q_+]\subset(0,1)$; when it is not a point, $\mu_{\beta,h}$ has atoms at both ends and a smooth positive density in between.
\item \parencite[Thm.~1.2 and Cor.~1.4]{AronowLopatto2026b} For every $h$, the law of the overlap array $(R_{ab})_{a<b}$ under $\E\gibbs\cdot$ converges to the overlap array of the Ruelle cascade of $\mu_{\beta,h}$ (with independent signs when $h=0$); in particular the law of $R_{12}$, of $\abs{R_{12}}$ when $h=0$, converges to $\mu_{\beta,h}$, and the limiting overlaps (in absolute value when $h=0$) are ultrametric.
\end{enumerate}
\end{theorem}

Before part~3, the overlap law and ultrametricity were known only for Gibbs measures that satisfy the Ghirlanda--Guerra identities \parencite{Panchenko2013a,Panchenko2013b}, such as those of generic mixed $p$-spin models or of perturbed Hamiltonians; SK itself is not generic. All three papers state that they were written with the assistance of large language models, used to suggest arguments, to draft and to write code for checks; for part~2 the observables and inequalities of the proof were located, in Lopatto's words, using \enquote{a massive number of numerical experiments orchestrated by large language models}. Building on an earlier version of part~1, \textcite{ChenHB2026} claims a proof that $q_\beta\to1$ as $\beta\to\infty$ and that the zero-temperature Parisi measure has support $[0,1)$; its arXiv record states that \enquote{ChatGPT 5.6 is the actual author}, and that Chen, listed to comply with arXiv policy, accepts full responsibility. At zero temperature \textcite{AuffingerChen2017} proved the Parisi formula for the ground-state energy, and \textcite{AuffingerChenZeng2020} that its Parisi measure has infinitely many points in its support.

\section{Measuring it}\label{sec:sk:measuring}

Two scripts do the measuring, in about one and two minutes of computation: \texttt{03-sk-replica.py} solves the replica equations quoted in \cref{sec:sk:physics}, and \texttt{03-sk-enum.py} enumerates all $2^N$ configurations for $N\le16$, so that each $\log Z_N$ and each Gibbs average is exact for its draw of the couplings, and averages over independent draws.\sidenote{A Gibbs measure costs about $2^NN^2$ operations, so $N=16$, with $65{,}536$ configurations, is the practical limit of exact enumeration on a laptop; Monte Carlo reaches larger $N$ but loses the exactness that makes the identity checkable at each size.}

\begin{figure}[t]
\begin{fullwidth}
\centering
\begin{tikzpicture}
\begin{groupplot}[group style={group size=3 by 1,horizontal sep=15mm},rplot,scale only axis,width=41mm,height=33mm]
\nextgroupplot[xlabel={inverse temperature $\beta$},ylabel={free entropy},xmin=0,xmax=2.05,ymin=0.66,ymax=1.68,
  xtick={0,0.5,1,1.5,2},title={(a) $\phi_{16}(\beta)$}]
\addplot[proofpred] table[x=beta,y=ann] {figures/data/03-sk-phi.dat};
\addplot[physpred] table[x=beta,y=phiRS] {figures/data/03-sk-phi.dat};
\addplot[physpred2] table[x=beta,y=phi1] {figures/data/03-sk-phi.dat};
\addplot[measured,ebars] table[x=beta,y=phi,y error plus=ci,y error minus=ci] {figures/data/03-sk-phi-enum.dat};
\nextgroupplot[xmode=log,ymode=log,xmin=3.6,xmax=15.5,xtick={4,6,8,10,12,14},xticklabels={4,6,8,10,12,14},
  log ticks with fixed point,xlabel={number of spins $N$},ymin=0.007,ymax=1.3,ylabel={$\skRS(q)-\phi_N$},
  title={(b) $\beta=0.8$, $q=0$},legend to name=skmeaslegend,legend columns=3,
  legend style={/tikz/every even column/.append style={column sep=4mm}}]
\addplot[proofpred] table[x=N,y=up8] {figures/data/03-sk-sumrule.dat};
\addplot[physpred] table[x=N,y=pred8] {figures/data/03-sk-sumrule.dat};
\addlegendimage{physpred2}
\addplot[measured,ebars] table[x=N,y=gap8,y error plus=gci8,y error minus=gci8] {figures/data/03-sk-sumrule.dat};
\addplot[measured2,ebars] table[x=N,y=rem8,y error plus=rci8,y error minus=rci8] {figures/data/03-sk-sumrule.dat};
\addplot[proofpred,forget plot] table[x=N,y=lo8] {figures/data/03-sk-sumrule.dat};
\legend{rigorous bound,replica prediction,one-step RSB,{$\phi_{16}$ in (a); the gap in (b), (c)},remainder}
\nextgroupplot[xmode=log,ymode=log,xmin=3.6,xmax=15.5,xtick={4,6,8,10,12,14},xticklabels={4,6,8,10,12,14},
  log ticks with fixed point,xlabel={number of spins $N$},ymin=0.0015,ymax=1.3,title={(c) $\beta=2$, $q=q^\star$}]
\addplot[proofpred] table[x=N,y=lb2] {figures/data/03-sk-sumrule.dat};
\addplot[measured,ebars] table[x=N,y=gap2,y error plus=gci2,y error minus=gci2] {figures/data/03-sk-sumrule.dat};
\addplot[measured2,ebars] table[x=N,y=rem2,y error plus=rci2,y error minus=rci2] {figures/data/03-sk-sumrule.dat};
\end{groupplot}
\end{tikzpicture}

\vspace{1mm}
\pgfplotslegendfromname{skmeaslegend}
\end{fullwidth}
\caption{Measurements at $h=0$, with Student-$t$ $95\%$ intervals. (a) The free entropy $\phi_{16}(\beta)$, the mean of $\frac1{16}\log Z_{16}$ over $120$ draws of the couplings (the same draws at every $\beta$, seed $3016$), against the replica-symmetric prediction at its minimizing $q$, the one-step value \eqref{eq:sk:p1rsb}, and the limit $\log2+\beta^2/4$ of the annealed bound \eqref{eq:sk:annealed}, which is the limit of $\phi_N$ for $\beta\le1$ (\cref{thm:sk:hightemp}). (b), (c) Guerra's sum rule \eqref{eq:sk:sumrule} on both sides: the gap $\skRS(q)-\phi_N$, from $\log Z_N$ ($8000$ draws at $N=4$ down to $1000$ at $N=14$), and the remainder, from the Gibbs measures of the interpolation at six Gauss--Legendre nodes in $t$ ($2000$ down to $100$ draws), with independent draws (seeds $30000+100N+k$ and $30050+100N+k$). In (b) the dashed lines are the bounds of \cref{thm:sk:hightemp} (above) and \eqref{eq:sk:annealed} (below), and the solid line is the asymptotic prediction $-\frac1{4N}\log(1-\beta^2)$; in (c), $q^\star=0.5304$ and the dashed line is $\skRS(q^\star)-\min\skPar_1=0.0022$, a lower bound at every $N$ by Guerra's broken bound.}\label{fig:sk:measure}
\end{figure}

In \cref{fig:sk:measure}(a) every point lies below the annealed curve, which for $\beta\le1$ is the limit, as \cref{thm:sk:guerratoninelli} requires; at $\beta=0.5$, $\phi_{16}=0.7501\pm0.0015$, against $0.7517$ for the bound at $N=16$. At $\beta=2$, $\phi_{16}=1.417\pm0.021$, while the replica-symmetric and one-step values are $1.5848$ and $1.5826$: the finite-size gap is about $77$ times the difference between the ans\"atze, and at this size the free entropy cannot tell them apart.

Panels (b) and (c) measure the identity \eqref{eq:sk:sumrule}, whose two sides come from different objects; they agree within their intervals at every $N$ and both temperatures: at $\beta=0.8$ and $N=10$, $0.02158\pm0.00101$ against $0.02158\pm0.00025$; at $\beta=2$ and $N=12$, $0.180\pm0.006$ against $0.183\pm0.012$. At $\beta=0.8$ the gap lies between the rigorous bounds and approaches the prediction $0.2554/N$, slowly: $N$ times the gap grows from $0.194$ at $N=4$ to $0.233$ at $N=14$, and $N\E\gibbs{R_{12}^2}$ rises from $1.44$ to $1.79\pm0.07$ toward its predicted limit $1/(1-\beta^2)=2.78$. At $\beta=2$ the gap decreases a little more slowly than $1/N$ ($N$ times the gap grows from $1.5$ to $2.4$), though in the limit it is at least $0.0022$, and positive by \cref{thm:sk:toninelli}. No measurement contradicts a theorem, and none can see the breaking of replica symmetry.

\section{The status board}\label{sec:sk:status}

The free entropy of SK is settled: the limit exists, the Parisi formula is a theorem with a unique minimizer, replica symmetry holds without a field exactly up to $\beta=1$, and it fails below the de~Almeida--Thouless line. Full replica-symmetry breaking is a refereed theorem just below the critical temperature at zero field. As of October 2026, replica symmetry above the line with a field, full replica-symmetry breaking in the rest of the low-temperature phase, and the identification of $x(q)$ with the law of the overlap of SK itself are proved in preprints of April to October 2026, prepared, as they state, with the assistance of large language models, and not yet refereed.

\begin{statusboard}
\statusrow{The free entropy has a limit and self-averages}{\proved}{\textcite{GuerraToninelli2002}; \cref{thm:sk:guerratoninelli}}
\statusrow{Replica symmetry at $h=0$, $\beta\le1$: $\phi=\log2+\beta^2/4$}{\proved}{\textcite{AizenmanLebowitzRuelle1987}; \cref{thm:sk:hightemp}}
\statusrow{Replica symmetry for $h>0$ on the whole region $\skAT\le1$}{\provedpreprint}{arXiv:2604.11921v2; with rate $1/N$ where $\skAT<1$, arXiv:2608.23413v2; AI-assisted (\cref{thm:sk:atline})}
\statusrow{Replica symmetry below the line}{\knownfalse}{\textcite{Toninelli2002}, all $\beta>1$ at $h=0$; \cref{thm:sk:toninelli,comp:sk:zeroT}}
\statusrow{The Parisi formula, with a unique minimizer}{\proved}{\textcite[Thm.~3]{Guerra2003}, \textcite{Talagrand2006}; uniqueness, \textcite{AuffingerChen2015a}}
\statusrow{Full RSB at $h=0$ for $\beta>1$ close to $1$: support $[0,q_\beta]$}{\proved}{\textcite[Thm.~1]{Zhou2026}}
\statusrow{Full RSB everywhere below the line: $h=0$ and $\beta>1$; $h>0$ and $\skAT>1$}{\provedpreprint}{arXiv:2607.11756v4; arXiv:2610.00127v1 with \textcite{Toninelli2002}; AI-assisted (\cref{thm:sk:overlap})}
\statusrow{Ultrametricity of the overlaps}{\provedunder}{the Ghirlanda--Guerra identities (generic mixed $p$-spin models, perturbed Hamiltonians) \parencite{Panchenko2013a}; SK itself, next row}
\statusrow{In SK, the law of $R_{12}$ tends to the Parisi measure; ultrametric}{\provedpreprint}{arXiv:2609.17442v1; AI-assisted (\cref{thm:sk:overlap})}
\statusrow{The ground-state energy per spin is the zero-temperature Parisi value}{\proved}{\textcite{AuffingerChen2017}; its value, $-0.76317$, is a numerical evaluation \parencite{OppermannSchmidtSherrington2007}}
\end{statusboard}

\section{Where the translation fails}\label{sec:sk:fails}

\paragraph{What was assumed.} Every proof here uses Gaussian couplings, through integration by parts, and a complete graph, through a covariance \eqref{eq:sk:cov} that depends on the configurations only through their overlap. Guerra's bounds use the convexity of $\xi$, which fails for odd mixtures, where \textcite{Panchenko2014} needed ultrametricity. The field is deterministic and uniform. And in the inference reading SK is the posterior under a wrong model (\cref{ch:nishimori} treats the right one).

\paragraph{Where the physics is wrong, or not yet right.} Replica symmetry is wrong below the line, as its negative entropy says. One step of breaking is not right either: its entropy is still negative at low temperature (\cref{fig:sk:phase}(b)), and at zero temperature no finite number of steps suffices \parencite{AuffingerChenZeng2020}. The local stability criterion does not characterize replica symmetry in general: \textcite{MourratSchertzer2026} constructed, in a preprint, mixed $p$-spin models in which the generalized de~Almeida--Thouless criterion fails. And simulations at accessible sizes do not test the breaking through the free entropy (\cref{sec:sk:measuring}).

\paragraph{Open problems.} The 2026 preprints of \cref{thm:sk:atline,thm:sk:overlap} await independent verification and refereeing. For $\beta>1$ and $h=0$, \textcite{AronowLopatto2026a} show, in a preprint also written with the assistance of large language models, that $\Var\log Z_N$ lies between orders $N^{4/15}$ and $N^{7/15}$; the exact exponent is open. And if the local criterion fails for some mixed $p$-spin models, as the preprint of \textcite{MourratSchertzer2026} shows, which condition characterizes their replica-symmetric phase?

\begin{exercises}
\item \consolidation (Moments.) (a) Prove the formula for $\E Z_N$ in \eqref{eq:sk:annealed}. (b) For $h\ne0$, show that $\E Z_N^2/(\E Z_N)^2\ge e^{-\beta^2/2}\exp(N\beta^2\tanh^4h/2)$, so that the argument of \cref{thm:sk:hightemp} fails with a field. \hint{under the tilted product law the overlap has mean $\tanh^2h$; use Jensen.}\label{ex:sk:moments}

\item \consolidation (The Parisi functional by hand.) Using the two rules after \eqref{eq:sk:pde}, show that $\skPar(\delta_q)=\skRS(q)$ and that $\skPar(m\delta_{q_0}+(1-m)\delta_{q_1})$ is \eqref{eq:sk:p1rsb}.\label{ex:sk:pde}

\item \consolidation (No field, high temperature.) Using $\log\cosh x\ge x^2/2-x^4/12$, show that $\skRS(\beta,0;q)-\skRS(\beta,0;0)\ge\frac{\beta^2q^2}4(1-\beta^2)$, so that $\inf_q\skRS(\beta,0;q)=\log2+\beta^2/4$ when $\beta\le1$.\label{ex:sk:qzero}

\item \consolidation (The annealed bound from the overlap.) Using \cref{lem:sk:derivative} and $\gibbs{R_{12}^2}\ge1/N$, rederive \eqref{eq:sk:annealed}.\label{ex:sk:derivative}

\item \byhand (The replica Hessian.) For a matrix indexed by pairs $\{a,b\}\subset\{1,\dots,n\}$, $n\ge4$, with entries $c_2$, $c_1$, $c_0$ for pairs that coincide, share one index, or are disjoint, show that its eigenvalues are $c_2-2c_1+c_0$ (multiplicity $n(n-3)/2$), $c_2+(n-4)c_1-(n-3)c_0$ (multiplicity $n-1$) and $c_2+2(n-2)c_1+\frac{(n-2)(n-3)}2c_0$ (multiplicity $1$).\label{ex:sk:replicon}

\item \byhand\hard (The line near the critical point.) Expanding the fixed-point equation to third order and the condition $\skAT=1$ to second order in $\epsilon=\beta-1$, show that on the line $q=\epsilon+O(\epsilon^3)$ and $h^2=\frac43\epsilon^3+O(\epsilon^4)$, so that $H^2\simeq\frac43(1-T)^3$. \hint{expand $\E\tanh^2x$ and $\E\sksech^4x$, $x=\beta\sqrt qZ+h$, in powers of $\beta^2q$ and $h^2$.} The script finds $(1-T)^3/H^2=0.7485$ at $H=10^{-4}$.\label{ex:sk:smallfield}

\item \byhand (A third-order transition.) Show that at $h=0$ the replica-symmetric order parameter is $q=\frac12(\beta^2-1)+O((\beta^2-1)^2)$ and that $\skRS(q)-\log2-\frac{\beta^2}4=-\frac1{48}(\beta^2-1)^3+O((\beta^2-1)^4)$ as $\beta\downarrow1$. \hint{$\E\log\cosh(aZ)=\frac{a^2}2-\frac{a^4}4+\frac{a^6}3+O(a^8)$.} At $\beta=1.005$ the ratio of the exact difference to this term is $0.991$.\label{ex:sk:critical}

\item \translate (The Gaussian integral around the saddle.) The physicist's Laplace method around $Q_{12}=0$ gives $\E Z_N^2/(\E Z_N)^2\simeq e^{-\beta^2/2}(1-\beta^2)^{-1/2}$. Prove this limit for $\beta<1$, by the central limit theorem and uniform integrability from \cref{lem:sk:secondmoment}, and show that the ratio tends to infinity for $\beta\ge1$.\label{ex:sk:hs}

\item \translate\hard (From the free entropy to the overlap.) (a) Show from \eqref{eq:sk:sumrule} that at $h=0$ and $\beta<1$, $\int_0^\beta b\,\E\gibbs{R_{12}^2}_b\dd b\to0$, with $\gibbs\cdot_b$ the Gibbs measure at inverse temperature $b$. (b) Let $h=0$, $0<\beta<1$, $c_\beta$ as in \cref{thm:sk:hightemp}, and $\delta_N=(\log2+\beta^2/4)/N+\beta\sqrt{\log(2/c_\beta)/N}$. Prove that $\E\gibbs{R_{12}^2}\le\frac2\beta\sqrt{\delta_N}$ for every $N$. \hint{$\phi_N$ is convex in $\beta$, so $\partial_\beta\phi_N(\beta)$ is at least the slope of the chord from $\beta'=\beta-2\sqrt{\delta_N}$; bound the chord with \cref{thm:sk:hightemp} and \eqref{eq:sk:annealed}, and use \cref{lem:sk:derivative}.}\label{ex:sk:overlapfromphi}

\item \measurement (Both sides of the line.) Modify \texttt{03-sk-enum.py} to compute $\skRS(q^\star)-\phi_N$ at $T=0.5$ with $H=0.5$ (below the line, $\skAT=1.063$) and $H=1$ (above it, $\skAT=0.640$), for $N=6,\dots,14$. Observe that $N$ times the gap is nearly constant in both cases (about $1.1$ and $0.5$), and explain.\label{ex:sk:fieldmeasure}

\item \research (A conditional second moment.) By \cref{ex:sk:moments}, the plain second moment fails when $h\ne0$. \textcite{BrenneckeYau2022} proved replica symmetry on the region $\beta^2\E\sksech^2(\beta\sqrt q\,Z+h)\le1$ by a conditional second-moment method. The preprints of \cref{thm:sk:atline} reach the whole region $\skAT\le1$ through the Parisi formula or through interpolation and cavity estimates. Can a conditional second moment reach it too?\label{ex:sk:conditional}
\end{exercises}

\notesandrefs
The physics runs from \textcite{SherringtonKirkpatrick1975} through \textcite{deAlmeidaThouless1978} to \textcite{Parisi1979b,Parisi1980b,Parisi1983} and the ultrametricity of \textcite{MezardEtAl1984a}. In 2021 Parisi received half of the Nobel Prize in Physics, for discovering the interplay of disorder and fluctuations in physical systems. The rigorous theory runs from \textcite{AizenmanLebowitzRuelle1987} through \textcite{GuerraToninelli2002}, \textcite{Guerra2003}, \textcite{Toninelli2002} and \textcite{AizenmanSimsStarr2003} to \textcite{Talagrand2006}, and on to \textcite{Panchenko2013a,Panchenko2014}; Talagrand received the Abel Prize in 2024. The monographs of \textcite{Talagrand2011a} and \textcite{Panchenko2013b} develop the theory, and \textcite{MontanariSen2024} teach the replica and interpolation computations side by side for readers without physics.

The presentation is the book's own in two places: the explicit constants of \cref{thm:sk:hightemp} and \cref{ex:sk:overlapfromphi}, which assemble classical second-moment, concentration and convexity arguments, and the derivation of Toninelli's second-order coefficient \eqref{eq:sk:toninelli} through the heat equation, which we checked numerically. The 2026 results are reported as their authors state them, in the arXiv versions current on 25 October 2026.

\paragraph{Translating the literature.} \textcite{Talagrand2011a} and \textcite{Lopatto2026a} use this chapter's Hamiltonian; Talagrand writes $p_N(\beta,h)$ for $\phi_N$, $\nu(f)$ for $\E\gibbs f$ and $R_{1,2}$ for $R_{12}$, and Lopatto writes $F_N$, $F$ and $\alpha(\beta,h)$ for $\phi_N$, $\phi$ and $\skAT$. \textcite{KusuokaNakajima2026b} and \textcite{Lopatto2026b} sum over all ordered pairs with the factor $\beta/\sqrt{2N}$, the same covariance up to a constant and the same $\beta$. \textcite{Panchenko2013b}, \textcite{Panchenko2014}, \textcite{AuffingerChen2015a} and \textcite{Zhou2026} use $\xi(q)=\beta^2q^2$, so their $\beta$ is our $\beta/\sqrt2$ (Zhou's critical point is $1/\sqrt2$); \textcite{AronowLopatto2026b} write $\xi(q)=\sum_pb_p^2q^p$, and SK is $b_2=\beta/\sqrt2$. \textcite{Guerra2003} writes the functional with $x(q)=\mu([0,q])$ and puts the field inside the energy, so his $h$ is our physical field $H=h/\beta$. Physicists write $f=-\phi/\beta$, $T=1/\beta$, $H=h/\beta$ and $J=1$, use the inverse $q(x)$ of $x(q)$, and call $P(q)=x'(q)$ the overlap distribution. \textcite{AuffingerChen2017} use $\xi(s)=s^2/2$ and a function $\gamma$ playing the role of $\beta x$ as $\beta\to\infty$.

\printchapterbib
\end{refsection}
